\documentclass[conference]{IEEEtran}
\IEEEoverridecommandlockouts

\usepackage[T1]{fontenc}
\usepackage[utf8]{inputenc}
\usepackage{cite}
\usepackage{amsmath,amssymb,amsfonts}
\usepackage{algorithmic}
\usepackage{graphicx}
\usepackage{textcomp}
\usepackage{xcolor}
\usepackage{listings}
\usepackage{booktabs}
\usepackage{hyperref}
\usepackage{float}
\usepackage{caption}
\usepackage{subcaption}

% ─── Code listing style ──────────────────────────────────────────────────────
\lstdefinestyle{cstyle}{
  language=C,
  basicstyle=\ttfamily\footnotesize,
  keywordstyle=\color{blue}\bfseries,
  commentstyle=\color{gray}\itshape,
  stringstyle=\color{red},
  numberstyle=\tiny\color{gray},
  numbers=left,
  stepnumber=1,
  frame=single,
  breaklines=true,
  captionpos=b,
  tabsize=4,
  showstringspaces=false
}

\begin{document}

% ─── TITLE ───────────────────────────────────────────────────────────────────
\title{Row-Major vs Column-Major Array Traversal in C:\\
An Experimental Analysis of Cache Miss Impact on CPU Performance}

\author{
  \IEEEauthorblockN{Thed Arthur TOKO DIKONGUE}
  \IEEEauthorblockA{
    \textit{Student at ISTY - Versailles Saint-Quentin-en-Yvelines University} \\
    Vélizy-Vllacoublay \\
    thdarthur@gmail.com
  }
}

\maketitle

% ─── ABSTRACT ────────────────────────────────────────────────────────────────
\begin{abstract}
This paper presents an experimental study on the performance impact of memory
access patterns in C programs, with a specific focus on two-dimensional array
traversal. We demonstrate, through reproducible benchmarks and hardware
performance counters, that iterating a matrix with the column index in the outer
loop---i.e., \texttt{tab[j][i]} with \texttt{j} as the outer loop---results in
significant performance degradation compared to row-major traversal
(\texttt{tab[i][j]}, \texttt{i} outer). Experiments were conducted on an AMD
Ryzen 7 3700U (Zen+, x86-64, Ubuntu Linux) using GCC at \texttt{-O2}
optimization level. Results from \texttt{perf stat} hardware counters show that
column-major traversal generates approximately \textbf{20$\times$ more cache
misses} than row-major traversal, with a cache miss rate of 29.65\% versus
7.60\%, translating into a \textbf{3.7$\times$ slowdown} on a $4096 \times
4096$ integer matrix (64~MB). Static and dynamic analysis using MAQAO OneView
further reveals that the column-major inner loop achieves only 33\%
vectorization coverage versus 100\% for the row-major case, and an Array Access
Efficiency of 51.1\% versus 100\%, confirming that strided memory access
prevents effective SIMD utilization. This work provides a concrete,
hardware-validated demonstration of spatial locality principles, intended as
both a pedagogical reference and a practical guide for performance-sensitive
C and C++ development.
\end{abstract}

\begin{IEEEkeywords}
cache miss, spatial locality, memory access patterns, row-major order,
column-major order, CPU performance, SIMD vectorization, AMD Zen+,
hardware performance counters, MAQAO
\end{IEEEkeywords}

% ─── 1. INTRODUCTION ─────────────────────────────────────────────────────────
\section{Introduction}

Modern processors achieve high throughput largely through a multi-level cache
hierarchy (L1, L2, L3) that bridges the latency gap between the CPU core and
main memory (DRAM). The fundamental unit of data transfer between memory and
cache is the \textit{cache line}---typically 64 bytes on x86-64
architectures---meaning that even a single-byte read triggers a 64-byte
transfer from the next memory level.

Two-dimensional arrays in C are laid out in memory using \textit{row-major
order}: elements of the same row are stored at consecutive memory addresses.
When a program accesses \texttt{tab[i][j]} with \texttt{j} varying fastest
(the inner loop), the CPU reads consecutive memory locations, exploiting the
spatial locality principle. Conversely, when \texttt{j} is the outer loop and
\texttt{i} varies fastest in the inner loop, each successive access is
separated by $N \times \texttt{sizeof(int)}$ bytes, potentially causing a
cache miss on every access once the matrix exceeds the cache size.

This phenomenon is well-known in the computer architecture
literature~\cite{hennessy2017}, yet concrete, hardware-validated demonstrations
combining timing benchmarks, hardware counters, and static analysis tools remain
scarce in accessible form. This paper fills that gap by providing:

\begin{itemize}
  \item A reproducible C benchmark measuring both access patterns across seven
        matrix sizes from $64 \times 64$ to $4096 \times 4096$;
  \item Hardware performance counter measurements via Linux \texttt{perf stat},
        quantifying cache misses at the L1 and LLC levels;
  \item Static and dynamic analysis via MAQAO OneView, exposing vectorization
        degradation and array access efficiency caused by strided memory access;
  \item A comparative analysis across two compiler optimization levels
        (\texttt{-O0} and \texttt{-O2}).
\end{itemize}

The remainder of this paper is organized as follows.
Section~\ref{sec:bg} presents theoretical background on cache hierarchy and
spatial locality. Section~\ref{sec:setup} describes the experimental setup.
Section~\ref{sec:results} presents timing results. Section~\ref{sec:perf}
analyzes hardware counters. Section~\ref{sec:maqao} details MAQAO findings.
Section~\ref{sec:discussion} discusses implications, and
Section~\ref{sec:conclusion} concludes.

% ─── 2. BACKGROUND ───────────────────────────────────────────────────────────
\section{Theoretical Background}
\label{sec:bg}

\subsection{CPU Cache Hierarchy}

Contemporary x86-64 processors implement a multi-level cache hierarchy to
mitigate the latency disparity between processing units and DRAM. On the AMD
Ryzen 7 3700U used in this study:

\begin{itemize}
  \item \textbf{L1 data cache}: 32~KB per core, latency $\approx$4 cycles
  \item \textbf{L1 instruction cache}: 64~KB per core
  \item \textbf{L2 cache}: 512~KB per core (private), latency $\approx$12 cycles
  \item \textbf{L3 cache (LLC)}: 4~MB shared, latency $\approx$40 cycles
  \item \textbf{Main memory (DRAM)}: 5795~MB, latency $\approx$200--300 cycles
\end{itemize}

The ratio of cache hit latency to DRAM latency is typically 50:1 to 100:1.
A program generating frequent DRAM accesses due to cache misses therefore
experiences severe throughput degradation regardless of computational complexity.

\subsection{Cache Lines and Spatial Locality}

When the CPU accesses a memory address not present in any cache level (a
\textit{cache miss}), it fetches an entire cache line of 64 bytes---equivalent
to 16 consecutive 32-bit integers---from the next level. This mechanism
exploits \textit{spatial locality}: the empirical observation that programs
tend to access memory addresses proximate to recently accessed
ones~\cite{denning2005}.

For a C array \texttt{int tab[N][N]}, the element \texttt{tab[i][j]} resides
at address:

\begin{equation}
  \text{addr}(\texttt{tab[i][j]}) = \text{base} + (i \cdot N + j) \cdot
  \texttt{sizeof(int)}
  \label{eq:addr}
\end{equation}

\subsection{Access Pattern Analysis}

\textbf{Row-major traversal} (\texttt{j} inner): consecutive values of
\texttt{j} produce addresses differing by exactly 4 bytes. Accessing
\texttt{tab[i][0]} through \texttt{tab[i][15]} loads 16 integers from a single
cache line after one miss---a theoretical miss rate of $1/16 = 6.25\%$.

\textbf{Column-major traversal} (\texttt{i} inner): consecutive values of
\texttt{i} produce addresses differing by $N \times 4$ bytes. For $N = 4096$,
this stride equals 16,384 bytes---256 cache lines apart. Each inner loop
iteration accesses a cache line unlikely to be resident, producing a miss on
virtually every access.

Table~\ref{tab:miss_theory} summarizes the theoretical miss rates.

\begin{table}[H]
\caption{Theoretical Cache Miss Rate Analysis ($N = 4096$, 32-bit integers)}
\label{tab:miss_theory}
\centering
\begin{tabular}{lcccc}
\toprule
\textbf{Pattern} & \textbf{Inner} & \textbf{Stride} &
\textbf{Miss/N} & \textbf{Miss rate} \\
\midrule
Row-major    & \texttt{j} & 4~B      & 1  & $\sim$6.25\% \\
Column-major & \texttt{i} & 16,384~B & 16 & $\sim$100\%  \\
\bottomrule
\end{tabular}
\end{table}

% ─── 3. EXPERIMENTAL SETUP ───────────────────────────────────────────────────
\section{Experimental Setup and Environment}
\label{sec:setup}

\subsection{Hardware Environment}

The experiments were conducted on a machine equipped with an AMD Ryzen 7 3700U
processor (Zen+ microarchitecture, 14nm, 4 physical cores / 8 threads via SMT).
Fig.~\ref{fig:lstopo} shows the full memory topology as reported by
\texttt{lstopo}.

\begin{figure}[H]
  \centering
  \includegraphics[width=0.85\columnwidth]{lstopo.png}
  \caption{Hardware memory topology (lstopo): AMD Ryzen 7 3700U, 4 cores,
           L1d 32~KB, L2 512~KB per core, L3 4~MB shared, 5795~MB RAM.}
  \label{fig:lstopo}
\end{figure}

Each core has a private 32~KB L1 data cache, a 64~KB L1 instruction cache, and
a private 512~KB L2 cache. All four cores share a single 4~MB L3 (LLC). The
cache line size is 64~bytes (\texttt{clflush size: 64} in
\texttt{/proc/cpuinfo}). Total system RAM is 5795~MB on a single NUMA node.

\begin{table}[H]
\caption{Experimental Environment Summary}
\label{tab:env}
\centering
\begin{tabular}{ll}
\toprule
\textbf{Component} & \textbf{Specification} \\
\midrule
CPU            & AMD Ryzen 7 3700U (Zen+, 4C/8T, 2.3~GHz) \\
L1d / L1i      & 32~KB / 64~KB per core \\
L2             & 512~KB per core (private) \\
L3 (LLC)       & 4~MB shared \\
Cache line     & 64~bytes \\
RAM            & 5795~MB (single NUMA node) \\
OS             & Ubuntu Linux (kernel 6.x) \\
Compiler       & GCC (\texttt{-O0} and \texttt{-O2}) \\
Profiler       & Linux \texttt{perf} v6.17.13 \\
Analyzer       & MAQAO OneView \\
Timing         & \texttt{clock\_gettime(CLOCK\_MONOTONIC)} \\
Runs / size    & 5 (mean reported) \\
\bottomrule
\end{tabular}
\end{table}

\subsection{Benchmark Implementation}

Two traversal functions were implemented to isolate the effect of loop ordering.
Both apply the \texttt{sin()} function to each element, preventing dead code
elimination while introducing a realistic floating-point workload.

\begin{lstlisting}[style=cstyle, caption={Row-major traversal (cache-friendly)},
label={lst:row}]
void row_comp(int *mat, double *res, int size) {
    for (int i = 0; i < size; i++)
        for (int j = 0; j < size; j++)
            res[i*size + j] = sin(mat[i*size + j]);
}
\end{lstlisting}

\begin{lstlisting}[style=cstyle,
caption={Column-major traversal (cache-unfriendly)}, label={lst:col}]
void column_comp(int *mat, double *res, int size) {
    for (int j = 0; j < size; j++)
        for (int i = 0; i < size; i++)
            res[i*size + j] = sin(mat[i*size + j]);
}
\end{lstlisting}

Matrices are heap-allocated via \texttt{malloc}. Values are initialized with
\texttt{rand() \% 10}. The only structural difference between the two functions
is the order of the outer and inner loop indices.

\subsection{Matrix Sizes Tested}

\begin{table}[H]
\caption{Matrix Sizes and Cache Pressure}
\label{tab:sizes}
\centering
\begin{tabular}{ccl}
\toprule
\textbf{N} & \textbf{Size} & \textbf{Cache pressure} \\
\midrule
64   & 16~KB  & Fits entirely in L1d     \\
128  & 62~KB  & Fits in L2               \\
256  & 256~KB & At L2 boundary           \\
512  & 1~MB   & Between L2 and L3        \\
1024 & 4~MB   & At LLC boundary          \\
2048 & 16~MB  & Exceeds LLC              \\
4096 & 64~MB  & Full RAM access required \\
\bottomrule
\end{tabular}
\end{table}

% ─── 4. BENCHMARK RESULTS ────────────────────────────────────────────────────
\section{Benchmark Results}
\label{sec:results}

\subsection{Execution Time}

Fig.~\ref{fig:comparison} shows mean execution time as a function of matrix
size for both access patterns and both optimization levels.

\begin{figure}[H]
  \centering
  \includegraphics[width=\columnwidth]{comparison_O0_O2.png}
  \caption{Execution time (ms) vs. matrix size $N$: row-major vs. column-major,
           \texttt{-O0} and \texttt{-O2} (AMD Ryzen 7 3700U, 5 runs each).}
  \label{fig:comparison}
\end{figure}

The divergence between access patterns begins at $N = 1024$ (4~MB, at the LLC
boundary) and increases monotonically. At $N = 4096$, the column-major pattern
is \textbf{3.1$\times$ slower} under \texttt{-O2} and \textbf{2.6$\times$
slower} under \texttt{-O0}. Fig.~\ref{fig:ratio} shows the slowdown ratio
explicitly.

\begin{figure}[H]
  \centering
  \includegraphics[width=\columnwidth]{ratio_col_over_row.png}
  \caption{Column-major slowdown ratio relative to row-major. The cache
           saturation threshold at $N = 1024$ (LLC = 4~MB) is clearly visible.}
  \label{fig:ratio}
\end{figure}

\subsection{Raw Timing Data}

\begin{table}[H]
\caption{Mean Execution Times, \texttt{-O2} (ms), AMD Ryzen 7 3700U}
\label{tab:results_o2}
\centering
\begin{tabular}{ccccc}
\toprule
\textbf{N} & \textbf{Size} & \textbf{Row (ms)} & \textbf{Col (ms)} &
\textbf{Ratio} \\
\midrule
64   & 16~KB  & 0.140  & 0.111  & 0.79 \\
128  & 62~KB  & 0.578  & 0.472  & 0.82 \\
256  & 256~KB & 2.121  & 1.728  & 0.81 \\
512  & 1~MB   & 8.368  & 7.017  & 0.84 \\
1024 & 4~MB   & 37.46  & 49.25  & 1.31 \\
2048 & 16~MB  & 138.6  & 235.9  & 1.70 \\
4096 & 64~MB  & 551.1  & 1701.2 & 3.09 \\
\bottomrule
\end{tabular}
\end{table}

% ─── 5. HARDWARE COUNTER ANALYSIS ────────────────────────────────────────────
\section{Hardware Performance Counter Analysis}
\label{sec:perf}

Linux \texttt{perf stat} was used to collect hardware performance counters for
both patterns on a $4096 \times 4096$ matrix under \texttt{-O2}. Separate
binaries were used per pattern to isolate readings.

\begin{table}[H]
\caption{\texttt{perf stat} Results, $N = 4096$, \texttt{-O2}}
\label{tab:perf}
\centering
\begin{tabular}{lcc}
\toprule
\textbf{Counter} & \textbf{row\_comp} & \textbf{column\_comp} \\
\midrule
Cache misses         & 1,904,790    & 37,279,590    \\
Cache references     & 25,057,659   & 125,749,082   \\
L1 load misses       & 9,290,356    & 40,823,961    \\
Miss rate            & 7.60\%       & 29.65\%       \\
Elapsed time         & 552~ms       & 2034~ms       \\
\midrule
\textbf{Miss ratio (col/row)} &
\multicolumn{2}{c}{\textbf{$\approx$20$\times$}} \\
\bottomrule
\end{tabular}
\end{table}

Column-major traversal generates \textbf{20$\times$ more cache misses} than
row-major traversal. The L1 load miss count increases by a factor of 4.4, and
the cache miss rate rises from 7.60\% to 29.65\%, confirming that the CPU
fetches data from higher memory levels on nearly one in three accesses. The
elapsed time ratio (3.7$\times$) is lower than the miss ratio (20$\times$)
because the \texttt{sin()} computation partially hides memory latency through
out-of-order execution in the row-major case.

% ─── 6. MAQAO ANALYSIS ───────────────────────────────────────────────────────
\section{MAQAO Static and Dynamic Analysis}
\label{sec:maqao}

\subsection{Methodology}

MAQAO OneView was used to perform combined static (CQA) and dynamic (LProf)
analysis~\cite{maqao}. Two separate runs were conducted---one per access
pattern---to obtain isolated per-function metrics.

\subsection{Global Metrics}

Fig.~\ref{fig:maqao_global_col} and Fig.~\ref{fig:maqao_global_row} show the
Global Metrics panels for \texttt{column\_comp} and \texttt{row\_comp}
respectively.

\begin{figure}[H]
  \centering
  \includegraphics[width=0.85\columnwidth]{maqao_global_col.png}
  \caption{MAQAO Global Metrics --- \texttt{column\_comp}:
           Total Time 11.24~s, Array Access Efficiency \textbf{51.1\%}.}
  \label{fig:maqao_global_col}
\end{figure}

\begin{figure}[H]
  \centering
  \includegraphics[width=0.85\columnwidth]{maqao_global_row.png}
  \caption{MAQAO Global Metrics --- \texttt{row\_comp}:
           Total Time 4.25~s, Array Access Efficiency \textbf{100\%}.}
  \label{fig:maqao_global_row}
\end{figure}

The \textit{Array Access Efficiency} metric directly quantifies memory access
quality: MAQAO reports \textbf{100\%} for \texttt{row\_comp} and \textbf{51.1\%}
for \texttt{column\_comp}, confirming that nearly half of all memory bandwidth
is wasted in the column-major case.

\subsection{Expert Summary --- Loop Analysis}

Fig.~\ref{fig:maqao_loops_col} and Fig.~\ref{fig:maqao_loops_row} show the
Expert Summary loop tables for each pattern.

\begin{figure}[H]
  \centering
  \includegraphics[width=\columnwidth]{maqao_loops_col.png}
  \caption{MAQAO Expert Summary --- \texttt{column\_comp} (Loop 3,
           benchmark.c:14--15): 36.84\% app. coverage, 33.33\% vectorization.}
  \label{fig:maqao_loops_col}
\end{figure}

\begin{figure}[H]
  \centering
  \includegraphics[width=\columnwidth]{maqao_loops_row.png}
  \caption{MAQAO Expert Summary --- \texttt{row\_comp} (Loop 2,
           benchmark.c:8--9): 2.24\% app. coverage, 100\% vectorization.}
  \label{fig:maqao_loops_row}
\end{figure}

\subsection{Quantitative Comparison}

\begin{table}[H]
\caption{MAQAO OneView Metrics Comparison, $N = 4096$, \texttt{-O2}}
\label{tab:maqao}
\centering
\begin{tabular}{lcc}
\toprule
\textbf{Metric} & \textbf{row\_comp} & \textbf{column\_comp} \\
\midrule
Coverage (\% app. time)         & 2.24\%       & 36.84\%       \\
Vectorization ratio (\%)        & 100\%        & 33.33\%       \\
Vector length use (\%)          & 41.67\%      & 29.17\%       \\
Array Access Efficiency (\%)    & \textbf{100} & \textbf{51.1} \\
Total time (s)                  & 4.25         & 11.24         \\
CQA speedup if fully vectorized & 2.93$\times$ & 7.71$\times$  \\
\bottomrule
\end{tabular}
\end{table}

The coverage metric is the most striking result: \texttt{column\_comp} consumes
\textbf{36.84\%} of total application time versus \textbf{2.24\%} for
\texttt{row\_comp}---a 16$\times$ difference directly attributable to cache
miss stalls. The vectorization ratio drops from 100\% to 33\%: GCC generates
packed AVX instructions (\texttt{VMOVUPD}, \texttt{VINSERTF128}) for the
row-major loop but cannot apply these to non-contiguous strided accesses. This
explains the potential speedup of 7.71$\times$ if fully vectorized---the
column-major loop leaves substantial SIMD performance unrealized.

% ─── 7. DISCUSSION ───────────────────────────────────────────────────────────
\section{Discussion}
\label{sec:discussion}

\subsection{Cache Saturation Threshold}

Results identify a clear \textit{cache saturation threshold} at $N = 1024$
(4~MB), coinciding exactly with the LLC size of the Ryzen 7 3700U. Below this
threshold, the Zen+ hardware prefetcher compensates for the strided access
pattern---it detects the regular stride and preloads cache lines proactively,
making column-major paradoxically faster. Above the LLC boundary, the prefetcher
cannot keep up with demand and cache miss penalties dominate.

\subsection{Compiler Behavior}

Even under \texttt{-O2}, GCC does not automatically perform loop interchange to
transform column-major access into row-major order. This transformation is only
legal when loop iterations are provably independent. The programmer bears
primary responsibility for access pattern design.

\subsection{Compounding Effect: Vectorization Loss}

Beyond cache misses, the MAQAO analysis reveals that the strided access pattern
also degrades SIMD vectorization: the vectorization ratio drops from 100\% to
33\%, and Array Access Efficiency from 100\% to 51.1\%. This compounds the
performance penalty beyond what cache miss counts alone would predict.

\subsection{Practical Implications}

\begin{itemize}
  \item Always iterate C arrays in row-major order (\texttt{j} as inner loop);
  \item When column access is required by algorithm logic, consider transposing
        the matrix before the hot loop;
  \item The penalty is negligible for matrices fitting in L2 but severe beyond
        the LLC boundary;
  \item Modern prefetchers (Zen+, Apple M-series) partially compensate for
        small working sets but cannot eliminate the penalty at scale.
\end{itemize}

% ─── 8. CONCLUSION ───────────────────────────────────────────────────────────
\section{Conclusion}
\label{sec:conclusion}

This paper provided a hardware-validated, end-to-end experimental demonstration
of cache miss impact induced by column-major array traversal in C. Using timing
benchmarks, Linux \texttt{perf} hardware counters, and MAQAO static/dynamic
analysis on an AMD Ryzen 7 3700U, we showed that:

\begin{itemize}
  \item Column-major traversal generates \textbf{20$\times$ more cache misses}
        than row-major for a $4096 \times 4096$ matrix;
  \item This translates into a \textbf{3.7$\times$ execution time slowdown};
  \item Vectorization ratio drops from 100\% to 33\%, and Array Access
        Efficiency from 100\% to 51.1\%, revealing a compounding degradation;
  \item A clear cache saturation threshold exists at the LLC boundary (4~MB),
        below which the hardware prefetcher masks the effect.
\end{itemize}

Future work could extend this study to multi-threaded contexts (OpenMP), GPU
memory coalescing (CUDA), tiled matrix traversal as mitigation, and comparisons
across multiple microarchitectures (Intel Alder Lake, ARM Cortex-A).

% ─── REFERENCES ──────────────────────────────────────────────────────────────
\begin{thebibliography}{00}

\bibitem{hennessy2017}
J.~L. Hennessy and D.~A. Patterson,
\textit{Computer Architecture: A Quantitative Approach}, 6th~ed.
Morgan Kaufmann, 2017.

\bibitem{denning2005}
P.~J. Denning,
``The locality principle,''
\textit{Communications of the ACM}, vol.~48, no.~7, pp.~19--24, Jul. 2005.

\bibitem{drepper2007}
U.~Drepper,
``What every programmer should know about memory,''
\textit{Red Hat, Inc.}, Tech. Rep., 2007. [Online]. Available:
\url{https://people.freebsd.org/~lstewart/articles/cpumemory.pdf}

\bibitem{maqao}
D.~Barthou \textit{et al.},
``MAQAO: Modular assembler quality analyzer and optimizer for Itanium-2,''
in \textit{Proc. Int. Conf. on High Performance Computing}, 2005.

\bibitem{fog2023}
A.~Fog,
``Optimizing software in C++: An optimization guide for Windows, Linux,
and Mac platforms,'' Tech. Rep., 2023. [Online]. Available:
\url{https://www.agner.org/optimize/}

\end{thebibliography}

\end{document}